\documentclass[a4paper,11pt]{article}
\usepackage[T1]{fontenc}
\usepackage[polish]{babel}
\usepackage{lmodern}
\usepackage[margin=2.7cm]{geometry}
\usepackage[hidelinks]{hyperref}
\title{Trzy rodzaje źródeł w krótkim tekście}
\author{Przykład dydaktyczny}
\date{}
\begin{document}
\maketitle
Podręcznik OpenStax opisuje model wahadła i jego założenia~\cite{moebs2016}.
Wykres można wykonać programowo; publikacja Huntera przedstawia
bibliotekę Matplotlib~\cite{hunter2007}. Do sprawdzania obliczeń w Pythonie
przydaje się również oficjalna dokumentacja modułu
\texttt{math}~\cite{pythonmath2026}. Każde źródło pełni inną rolę:
podręcznik podaje teorię, artykuł opisuje narzędzie, a dokumentacja
określa działanie funkcji. Sam wykres z danych przykładowych nie staje
się dzięki tym cytowaniom wynikiem rzeczywistego pomiaru.
\bibliographystyle{plain}
\bibliography{references}
\end{document}
